% Source: https://artofproblemsolving.com/community/c3h1982197_problem_i_aint_the_smartest_child
% Stable identity: aops-topic-1982197-q01
Note that $(1+\sqrt3)^2=1+2\sqrt3+3=4+2\sqrt3$. Since $1+\sqrt3>0$, it is the (nonnegative) square root: $\sqrt{4+2\sqrt3}=1+\sqrt3$. Hence
\[
\sqrt{4+2\sqrt3}-\sqrt3=1+\sqrt3-\sqrt3=\boxed{1}.
\]
To find such a form in general, look for $\sqrt{a+2\sqrt b}=\sqrt x+\sqrt y$ with $x+y=a$ and $xy=b$; here $x+y=4$, $xy=3$ gives $x=3$, $y=1$.
